\chapter{Conclusion and Future Work}
\label{chapter-11-conclusion-and-future-work}

This chapter closes the thesis in three parts: a summary of the research methodology and the path from problem statement to the three contributions; a conclusion that directly answers the two research questions and positions the proposed architecture against the state of the art; and recommendations for future work organized into payload, communications, and mission-simulation themes.

%------------------------------------------------------------------
\section{Summary of Research}\label{sec:conclusion-summary}
%------------------------------------------------------------------

Wildfire risk in Mediterranean Europe --- particularly in Portugal --- has intensified, with longer fire seasons, larger burned areas, and extreme events that overwhelm tactical decision windows. Current satellite services face a resolution--revisit--latency gap: geostationary platforms monitor continuously but at coarse resolution, while polar LEO missions offer finer imagery at revisit intervals of hours to days. This thesis asked: \emph{What fire observation satellite constellation design can close the spatiotemporal gap, achieving roughly 30-minute revisit and approximately 6~m ground resolution over continental Portugal?} A secondary question is: \emph{What coverage does the proposed constellation provide over global fire-prone regions?}

The study began with a characterisation of the resolution--revisit--latency gap in current satellite fire-detection (Chapter~\ref{c1_introduction}), followed by the definition of mission requirements and the identification of global fire-prone domains from satellite-derived active-fire records (Chapter~\ref{chapter-2-the-wildfire-challenge-and-study-case-definition-from-global-threat-to-local-data-needs}).  A Repeating Ground-Track Orbit resonance analysis, generalised to arbitrary target latitudes, produced a family of RGTO modes for continental Portugal; the selected mode was embedded in a Walker constellation framework, and a parametric study of the sensor off-nadir angle determined the minimum spacecraft count for the required revisit cadence (Chapter~\ref{73-constellation-architectures-design}).  High-fidelity GMAT propagation validated orbital stability, and an observational-link geometry model quantified temporal coverage, revisit statistics, and ground sampling distance over both the Portuguese core domain and a global fire-prone domain (Chapter~\ref{chapter-9-the-proposed-system-concept}).  In parallel, a temporal-coincidence analysis of geostationary and multispectral LEO fire imagery demonstrated the coarse-to-fine tip-and-cue detection chain empirically (Chapter~\ref{sec:proposed_architecture}).  These elements converged into three contributions: an RGTO constellation tuned for Portugal, a two-layer \emph{Watcher-Zoom} architecture coupling existing geostationary fire-watch assets with an agile LEO segment, and a quantitative assessment of the constellation's natural extension to global fire-prone regions.



%------------------------------------------------------------------
\section{Conclusion}\label{sec:conclusion-answer}
%------------------------------------------------------------------

The constellation $\mathcal{C}_{\text{Zoom}} = \mathcal{C}_{15/5/72}$---fifteen satellites in five planes at approximately 488~km, inclination matched to Portugal---delivers an average revisit of roughly 21~minutes (worst-case about 35~minutes) and a ground sampling distance below 6~m per pixel, satisfying requirements \reqRef{R7} and \reqRef{R8}.  The orbit resonance concentrates passes over the target latitude, yielding roughly twice the daily coverage of a comparable Sun-Synchronous Orbit with half the satellite count.  The \emph{Watcher-Zoom} tip-and-cue layer complements this design: existing geostationary fire-watch platforms (GOES, Meteosat, Himawari) provide coarse alerts that cue the LEO satellites for rapid, narrow-field imaging of confirmed hotspots, so each satellite transmits only fire-event data rather than full-swath imagery.   A demonstration of the Watcher-Zoom concept was conducted by testing coarse-to-fine chain image alignment of fire detections using geostationary coarse but fast images from FIRMS and Sentinel-2's slower images with high spatial resolution, confirming the concept.

Designed for Portugal at the minimum viable size, the coverage footprint extends naturally across the globe.  The temperate belt — southern Europe, the western United States, southeastern Australia, southern Africa — receives 6 to 7 hours of daily coverage with worst-case revisit gaps of about 30 minutes.  At the equatorial fringe---Amazonia, sub-Saharan Africa, maritime South-East Asia---coverage persists above 3.5 hours per day with an average revisit around 40 to 50 minutes, still outperforming current polar-orbiting fire sensors at comparable resolution.  A constellation sized for one country, therefore, services the majority of fire-prone nations, turning a regional mission into a global one.

No known Earth-observation constellation reviewed in this study achieves a fire-detection performance combining sub-30-minute revisit with metre-scale thermal resolution over Portugal or any other fire-prone region examined herein.  By uniting a resonant orbit, a geostationary-cued tip-and-cue layer, and their emergent global reach, the proposed architecture closes this gap from just fifteen satellites---offering regions facing the same wildfire threat a platform for cooperative observation.  This work lays the scientific foundation for a next-generation \emph{Satellite Constellation for Fast Fire Detection: Mission Design for Portugal Covering Global Fire-Prone Regions}.

%------------------------------------------------------------------
\section{Recommendations for Future Work}\label{sec:conclusion-future}
%------------------------------------------------------------------

The following recommendations trace a path from component-level design to full mission validation.

\subsection{Payload and Imaging}

\begin{enumerate}
    \item \textbf{End-to-end image simulator and payload constraints.} The FIRMS active-fire dataset, Sentinel-2 fire-image archive, and orbital simulation provide the inputs for an end-to-end image simulator of the Watcher-Zoom concept, enabling quantitative derivation of spectral-band selection, radiometric sensitivity, pointing-agility requirements, and the communications budget, bridging this scientific study into the engineering domain.
    \item \textbf{Image-quality simulation.} Model the imaging chain---motion blur, jitter, and off-nadir degradation---to determine the maximum usable field of view and achievable resolution under realistic attitude-control performance.
    \item \textbf{Co-design with machine vision.} Investigate co-design of the optical payload with onboard machine-learning models; a lower raw resolution may be acceptable if a trained model can reliably detect fire signatures and smoke plumes on the sensor output, reducing hardware complexity and downlink volume.
\end{enumerate}

\subsection{Communications and Autonomy}

\begin{enumerate}
    \item \textbf{Ground-segment and relay architecture.} Design the end-to-end communications path from the Zoom layer to fire-management end users, comparing direct ground-station contact, inter-satellite optical links, and relay through commercial LEO broadband constellations.
    \item \textbf{Onboard autonomy.} Develop autonomous operations concepts---onboard fire detection, dynamic re-tasking of the pointing system, and autonomous channel acquisition---to minimise ground-in-the-loop latency.
\end{enumerate}

\subsection{Mission Simulation and Scaling}

\begin{enumerate}
    \item \textbf{Full mission scenario.} Extend the geometric-access simulation to encompass sensor scheduling, detection probability, relay-chain data propagation, and integration with ground-based fire-management systems.
    \item \textbf{Station-keeping and propulsion.} Size the propulsion subsystem required to maintain RGTO resonance over the mission lifetime, and perform conjunction-screening and collision-avoidance analysis.
    \item \textbf{Global extension and cost--benefit analysis.} Evaluate constellation performance over the full global fire-prone domain $\mathcal{D}_{\mathrm{global}}$ and develop a cost--benefit model quantifying the socio-economic impact of reduced fire-detection latency.
\end{enumerate}